\documentclass[11pt]{article}
\usepackage{acl}
\usepackage{times}
\usepackage{latexsym}
\usepackage[T1]{fontenc}
\usepackage[utf8]{inputenc}
\usepackage{microtype}
\usepackage{inconsolata}
\usepackage{graphicx}
\usepackage{booktabs}
\usepackage{multirow}

\title{Text or Image? How Task Type Determines Modality Dominance in Vision-Language Models}

\author{
Rodela Ghosh \\
University of South Florida \\
\texttt{rg21@usf.edu} \\
\And
Aviral Gupta \\
University of South Florida \\
\texttt{aviralgupta@usf.edu} \\
}

\begin{document}
\maketitle

\begin{abstract}
Vision-language models (VLMs) can process both text and images, yet it remains unclear which modality dominates when the two conflict, and whether this depends on task type. We evaluate eight open-source VLMs on GSM8K---six of them across a full eight-benchmark suite---spanning two protocols: Protocol~A, where text problems are rendered as images (GSM8K, SVAMP, MATH-500, AQuA-RAT), and Protocol~B, where vision provides unique information unavailable in text (MathVista, AI2D, ChartQA, ScienceQA). We introduce a novel \textit{mismatch} condition in which the image and text present different problems, directly measuring modality preference under conflict. Under Protocol~A, rendering problems as images reduces accuracy in most models (drops of 3--54pp), and nearly all models show overwhelming text preference in the mismatch condition (87--100\%). Under Protocol~B, the pattern reverses: image conditions outperform text-only by 10--54pp across all models and benchmarks. Our results reveal that modality dominance is task-conditional: VLMs default to text when vision is redundant, but correctly leverage visual input when it provides unique information. This suggests current VLMs are not text-first by design, but by the nature of what each modality offers.
\end{abstract}

\section{Introduction}

Vision-language models (VLMs) have demonstrated strong capabilities across tasks that combine visual perception and language understanding~\cite{liu2023llava,bai2023qwen,chen2024internvl2}. Mathematical reasoning is a particularly demanding test of these capabilities, requiring precise symbol recognition, multi-step inference, and reliable numerical grounding~\cite{lu2023mathvista,cobbe2021gsm8k}.

Recent work has documented a consistent \textit{modality gap}: VLMs tend to perform worse when mathematical problems are presented as images rather than text~\cite{zhang2025gsm8kv,wang2026crossmath,chen2026vista}. This gap has been attributed primarily to perceptual failures at the image-reading stage rather than downstream reasoning breakdowns~\cite{liu2026readingnotthinking}, and is more pronounced in smaller models~\cite{zhang2026vistira}. Despite this growing body of work, two questions remain unanswered: \textit{(1) when a model receives conflicting information across modalities, which modality does it follow?} and \textit{(2) does this preference depend on the nature of the task?}

Prior studies vary the modality of a single input but do not pit the two modalities against each other, nor do they systematically compare tasks where vision is redundant versus informative. We address both gaps with a two-protocol evaluation design. Under \textbf{Protocol~A} (rendered-text benchmarks), the visual input is a rendered image of the text problem---vision is redundant. Under \textbf{Protocol~B} (natural visual benchmarks), the image provides information that cannot be recovered from text alone---vision is essential. We additionally introduce a \textbf{mismatch condition} in Protocol~A that directly measures modality preference by presenting conflicting image and text inputs.

We evaluate eight open-source VLMs (2B--8B parameters) on GSM8K, six of them across all eight benchmarks. Our contributions are:

\begin{itemize}
    \item \textbf{Novel mismatch condition:} the first evaluation design that directly measures modality preference by forcing VLMs to choose between conflicting visual and textual inputs.
    \item \textbf{Task-conditional dominance:} text dominates when vision is redundant (Protocol~A: 87--100\% text preference); image dominates when vision is informative (Protocol~B: +10--54pp image advantage).
    \item \textbf{Reasoning-trace validation:} a post-hoc rescore of ambiguous mismatch trials confirms text dominance is genuine, not an artifact of answer overlap.
    \item \textbf{Architectural divide:} across Protocol~A benchmarks, we identify resilient models (near-zero accuracy drop) and vulnerable models (14--54pp drop), and characterize what drives this split.
\end{itemize}

Our central finding is that VLMs are not unconditionally text-first: they follow the modality that provides the most useful information for the task at hand.

\section{Related Work}

\paragraph{Modality gap in VLMs.}
Several recent studies have shown that rendering text as images degrades VLM performance on reasoning tasks. \citet{zhang2025gsm8kv} converted GSM8K problems into comic-style visual scenes and found large drops in even strong closed-source models (Gemini-2.5-Pro: 95\%$\to$47\%). \citet{wang2026crossmath} introduced CrossMath, a benchmark with text-only, image-only, and image+text formats, and found VLMs reason primarily in textual space. \citet{chen2026vista} benchmarked 30+ VLMs on visualized text and found over 70\% of errors are perceptual rather than reasoning failures. \citet{liu2026readingnotthinking} specifically studied whether the gap stems from reading failures or reasoning failures, finding modality gaps of up to 40.6pp. Our work extends these studies by introducing a mismatch condition and contrasting rendered-text tasks with genuinely visual tasks.

\paragraph{VLMs for mathematical and visual reasoning.}
Benchmarks such as MathVista~\cite{lu2023mathvista}, MATH~\cite{hendrycks2021math}, and VisAidMath~\cite{wu2024visaidmath} evaluate VLMs on problems in their natural format. AI2D~\cite{kembhavi2016ai2d}, ChartQA~\cite{masry2022chartqa}, and ScienceQA~\cite{lu2022scienceqa} require genuine visual understanding. Our Protocol~B uses these benchmarks to show that VLMs do leverage vision when it is informative.

\paragraph{Modality preference and text bias.}
\citet{ijcai2025wordspixels} showed through a broad survey that VLMs frequently rely on textual cues over visual information and exhibit weak spatial reasoning. \citet{zhang2026vistira} proposed closing the modality gap via structured tool integration. Our results qualify this text-bias claim: it holds specifically when vision is redundant, not universally.

\paragraph{VLM robustness.}
\citet{chen2026vlmrobust} evaluated VLM robustness under realistic corruptions, finding that text recognition degrades most under blur and noise. This motivates noise ablation experiments we leave for future work.

\section{Methodology}

\subsection{Evaluation Protocols}

We evaluate VLMs under two protocols:

\textbf{Protocol~A (rendered-text benchmarks):} The problem text is rendered as a 900px-wide PNG image (DejaVu Sans, black on white). Three conditions are run: text-only (C1), rendered image (C2), and mismatch (C3). Vision is redundant --- the image contains no information beyond what is in the text prompt.

\textbf{Protocol~B (natural visual benchmarks):} Problems are presented in their original multimodal format. Two conditions are run: text-only (the question text without the image) and multimodal (question text plus the original image). Vision is essential --- the image provides diagrams, charts, or figures required to answer correctly.

\subsection{Protocol A: Conditions and Scoring}

Under Protocol~A, three conditions are evaluated on each benchmark:

\begin{itemize}
    \item \textbf{C1 (Text-only):} Plain text prompt with the instruction ``Solve the following math problem step by step. End with `\#\#\#\# $\langle$answer$\rangle$'.'' No image provided.
    \item \textbf{C2 (Rendered image):} Model receives only the rendered PNG; no problem text in the text channel.
    \item \textbf{C3 (Mismatch):} Model receives image $i$ (problem $i$) paired with text prompt for problem $i{+}1$. Problems have distinct answers in 98.9\% of cases; equal-answer pairs are excluded as ambiguous.
\end{itemize}

For C3, we classify responses as \textit{image}, \textit{text}, \textit{neither}, \textit{ambiguous}, or \textit{invalid} by exact numeric match. We report text preference over decidable trials only. We additionally perform a \textbf{reasoning-trace rescore} on \textit{neither} trials: unique keywords and numbers from each problem are matched against the model's chain-of-thought, with keyword hits weighted 2$\times$ over number hits. Trials are reclassified as \textit{text\_reasoning} or \textit{image\_reasoning} if one problem scores strictly higher.

\subsection{Benchmarks}

\noindent\textbf{Protocol~A:} GSM8K~\cite{cobbe2021gsm8k} ($N{=}1{,}319$), SVAMP ($N{=}300$), MATH-500~\cite{hendrycks2021math} ($N{=}500$), AQuA-RAT ($N{=}254$). All rendered datasets are publicly available.\footnote{\url{https://huggingface.co/vlm-modality-research}}

\noindent\textbf{Protocol~B:} MathVista~\cite{lu2023mathvista} ($N{=}1{,}000$), AI2D ($N{=}1{,}000$), ChartQA ($N{=}1{,}000$), ScienceQA ($N{=}1{,}000$).

\subsection{Models}

We evaluate eight open-source VLMs (Table~\ref{tab:models}), spanning 2B--8B parameters and six architectural families. All eight are evaluated on GSM8K; six (excluding Qwen2.5-VL-7B and Phi-3.5-Vision) are run across the full eight-benchmark suite. All models are run in \texttt{bfloat16} without quantization, using greedy decoding, on single-GPU nodes.

\begin{table}[h]
\centering
\small
\begin{tabular}{lll}
\toprule
\textbf{Model} & \textbf{Params} & \textbf{Family} \\
\midrule
Qwen2-VL-2B-Instruct        & 2B  & Qwen    \\
LLaVA-1.6-Mistral-7B        & 7B  & LLaVA   \\
Qwen2.5-VL-7B-Instruct      & 7B  & Qwen    \\
Idefics3-8B-Llama3          & 8B  & Idefics \\
MiniCPM-V-2.6               & 8B  & MiniCPM \\
InternVL2-8B                & 8B  & InternVL\\
LLaVA-OneVision-Qwen2-7B    & 7B  & LLaVA   \\
Phi-3.5-Vision-Instruct     & 4B  & Phi     \\
\bottomrule
\end{tabular}
\caption{Models evaluated in this study. All eight are evaluated on GSM8K; the
first six are additionally evaluated across the full eight-benchmark suite.}
\label{tab:models}
\end{table}

\subsection{Statistical Analysis}

We report accuracy with 95\% Clopper-Pearson confidence intervals. We use McNemar's test to assess whether accuracy differences between conditions are statistically significant, and Cohen's $h$ as an effect size measure. For mismatch preference, we apply a chi-squared test on decidable trials. Code is available at \url{https://github.com/Ro-netizen004/vlm-modality-investigation}.

\section{Results}

\subsection{Protocol A: Rendered-Text Benchmarks}

\paragraph{Accuracy (C1 vs.\ C2).}
Table~\ref{tab:accuracy_gsm8k} shows GSM8K results for all eight models. Two groups emerge: \textbf{resilient models} (Qwen2.5-VL-7B: $+$0.2pp; InternVL2-8B: $+$0.5pp; Qwen2-VL-2B: $-$2.6pp, all ns) and \textbf{vulnerable models} (Idefics3: $-$49.3pp, $h{=}1.03$; LLaVA-OneVision: $-$16.4pp; MiniCPM: $-$14.7pp; LLaVA-1.6: $-$14.3pp; Phi-3.5: $-$4.9pp; all $p{<}0.001$). This pattern holds across SVAMP and MATH-500 for the six models with full-suite coverage, with Idefics3 consistently showing the largest drops (SVAMP: $-$53.7pp; MATH-500: $-$22.6pp).

\begin{table}[h]
\centering
\small
\begin{tabular}{lcccc}
\toprule
\textbf{Model} & \textbf{C1} & \textbf{C2} & \textbf{Drop} & \textbf{$h$} \\
\midrule
Qwen2-VL-2B        & .426 & .400 & $-$.026\rlap{$^\dagger$} & .052 \\
LLaVA-1.6          & .434 & .291 & $-$.143*** & .299 \\
Qwen2.5-VL-7B      & .845 & .847 & $+$.002\rlap{$^\dagger$} & .004 \\
Idefics3-8B        & .773 & .280 & $-$.493*** & 1.033 \\
MiniCPM-V-2.6      & .748 & .601 & $-$.147*** & .316 \\
InternVL2-8B       & .771 & .776 & $+$.005\rlap{$^\dagger$} & .013 \\
LLaVA-OneVision    & .790 & .626 & $-$.164*** & .363 \\
Phi-3.5-Vision     & .776 & .727 & $-$.049*** & .112 \\
\bottomrule
\end{tabular}
\caption{GSM8K accuracy under C1 (text-only) and C2 (rendered image). $^\dagger$ns; ***$p{<}0.001$ (McNemar). $h$ = Cohen's $h$.}
\label{tab:accuracy_gsm8k}
\end{table}

\paragraph{Mismatch condition (C3).}
Table~\ref{tab:mismatch} shows that, with one exception, all models demonstrate near-unanimous text preference in the mismatch condition, ranging from 96.2\% to 100\% on GSM8K after reasoning-trace rescore, with true ``neither'' responses accounting for at most 2\% of trials. The exception is \textbf{Phi-3.5-Vision}, whose text preference falls to 86.9\% with 5.8\% true-neither: nearly a quarter of its reasoning-rescored trials engage the image problem, making it the only model that attends to the visual input a non-trivial fraction of the time. The dominant pattern holds across SVAMP (96--100\%) and MATH-500 (79--93\%) for the full-suite models, with the lower MATH-500 values reflecting increased difficulty and ambiguity of symbolic answers.

\begin{table}[h]
\centering
\small
\begin{tabular}{lccc}
\toprule
\textbf{Model} & \textbf{TextPref} & \textbf{TextPref$^*$} & \textbf{Neither$^*$} \\
\midrule
Qwen2-VL-2B        & .983 & .988 & 0.2\% \\
LLaVA-1.6          & .989 & .994 & 0.0\% \\
Qwen2.5-VL-7B      & .993 & .994 & 0.0\% \\
Idefics3-8B        & .994 & .994 & 2.0\% \\
MiniCPM-V-2.6      & .994 & .995 & 0.0\% \\
InternVL2-8B       & .962 & .962 & 0.1\% \\
LLaVA-OneVision    & .998 & .998 & 0.4\% \\
Phi-3.5-Vision     & .987 & .869 & 5.8\% \\
\bottomrule
\end{tabular}
\caption{GSM8K mismatch modality preference. TextPref = exact-match text preference over decidable trials. TextPref$^*$, Neither$^*$ = after reasoning-trace rescore.}
\label{tab:mismatch}
\end{table}

\subsection{Protocol B: Natural Visual Benchmarks}

Table~\ref{tab:protocolb} shows results under Protocol~B. The pattern reverses completely: image conditions outperform text-only across every model and benchmark. The gains are largest on ChartQA (+29--54pp), where text-only accuracy is near zero (1.5--4\%) since chart questions cannot be answered without the visual. AI2D (+10--23pp) and ScienceQA (+11--30pp) show consistent gains. MathVista, which combines visual and mathematical reasoning, shows moderate but significant image advantage (+14--21pp).

\begin{table}[h]
\centering
\small
\setlength{\tabcolsep}{4pt}
\begin{tabular}{lcccc}
\toprule
\textbf{Model} & \textbf{MathV.} & \textbf{AI2D} & \textbf{ChartQA} & \textbf{SciQA} \\
\midrule
Qwen2-VL-2B   & $+$.146 & $+$.158 & $+$.514 & $+$.157 \\
LLaVA-1.6     & $+$.046 & $+$.099 & $+$.288 & $+$.109 \\
Idefics3-8B   & $+$.159 & $+$.158 & $+$.464 & $+$.218 \\
MiniCPM-V-2.6 & $+$.208 & $+$.116 & $+$.540 & $+$.116 \\
InternVL2-8B  & $+$.176 & $+$.234 & $+$.406 & $+$.298 \\
LLaVA-OV      & $+$.136 & $+$.171 & $+$.539 & $+$.240 \\
\bottomrule
\end{tabular}
\caption{Protocol~B image advantage (image accuracy $-$ text-only accuracy) per benchmark. All differences are significant ($p{<}0.001$, McNemar's test). MathV.\ = MathVista; SciQA = ScienceQA; LLaVA-OV = LLaVA-OneVision.}
\label{tab:protocolb}
\end{table}

\subsection{Noise Robustness}
\label{sec:noise}
% PHASE 4 — drop in noise ablation results here.
% Expected story: rendered-image accuracy (Protocol A) degrades further under
% Gaussian/blur/salt-and-pepper noise, confirming that vision is an unreliable
% channel for redundant text. Compare degradation slopes across resilient vs.
% vulnerable models.
\textit{[Phase 4 results pending.]} We evaluate the stability of the rendered-image condition under increasing image corruption (Gaussian noise, blur). This isolates whether the Protocol~A gap is driven by fragile perception.

\subsection{Prompt Sensitivity}
\label{sec:prompt}
% PHASE 5 — drop in prompt-sensitivity results here.
% Expected story: can prompting shift mismatch modality preference? Test prompts
% that explicitly instruct the model to attend to the image. Report whether text
% dominance is robust (unchanged) or steerable (preference shifts).
\textit{[Phase 5 results pending.]} We test whether the text dominance observed in the mismatch condition can be overridden through prompts that explicitly direct attention to the image, quantifying how robust the preference is to instruction.

\subsection{Mechanistic Analysis}
\label{sec:mechanistic}
% PHASE 6 — drop in attention/probing results here.
% Expected story: in the mismatch condition, attention mass concentrates on text
% tokens over image tokens, providing a mechanistic account of text dominance.
% Keep claims modest if results are noisy — frame as "preliminary evidence."
\textit{[Phase 6 results pending.]} To explain \textit{why} text dominates under conflict, we examine the distribution of attention across text and image tokens in the mismatch condition, providing preliminary mechanistic evidence for the observed preference.

\section{Analysis}

\subsection{Task-Conditional Modality Dominance}

The contrast between Protocol~A and Protocol~B reveals a principled pattern: modality dominance tracks information utility. Under Protocol~A, the image is a rendering of the text---it adds no new information---and models default to the more reliable text channel. Under Protocol~B, the image contains information absent from the text (charts, diagrams, figures), and models correctly leverage it. This is not a universal text bias; it is a \textit{redundancy-driven} text preference.

Notably, this holds even for Idefics3, which performs poorly on Protocol~A rendered images (suggesting weak OCR) but shows strong image advantage on Protocol~B tasks (ChartQA: $+$46.4pp, ScienceQA: $+$21.8pp). The model can leverage visual information when it is genuinely visual---its failure under Protocol~A is specifically an OCR/text-reading failure.

\subsection{Disagreement Analysis}

Among Protocol~A vulnerable models, disagreements between C1 and C2 strongly favor text (e.g., 691 text-advantage vs.\ 41 image-advantage for Idefics3 on GSM8K). For resilient models, disagreements are balanced (150 vs.\ 157 for InternVL2), consistent with random variation. Problems where both conditions fail are consistently longer, contain more numbers, and have larger answer magnitudes---suggesting intrinsic difficulty rather than modality effects. For Idefics3, question length negatively correlates with the modality gap ($r = -0.20$, $p < 10^{-12}$): shorter, simpler problems are where vision most hurts the model.

\subsection{Why Are Some Models Resilient?}

Resilient models (InternVL2-8B) are also the strongest performers on text-only tasks. We hypothesize that stronger OCR capability and more diverse multimodal pretraining allow these models to read rendered text reliably. This is consistent with Idefics3's high text-only accuracy (77.3\%) but near-chance rendered-image accuracy (28.0\%)---the reasoning capability is intact, but the image-reading step fails.

\section{Limitations}

Our study has several limitations. First, the mismatch condition uses numeric-answer scoring and therefore does not extend cleanly to multiple-choice benchmarks (AQuA-RAT), where letter answers prevent reliable image-vs-text attribution. Second, all evaluated models are open-source and in the 2B--8B range; we do not test frontier closed-source models (e.g., GPT-4o, Gemini), so generalization to larger models is unverified. Third, the reasoning-trace rescore relies on lexical keyword and number overlap, which may misattribute trials where a model paraphrases problem content. Fourth, the mismatch condition is currently defined only for Protocol~A; we do not yet probe modality conflict in genuinely visual tasks. Finally, our rendering pipeline uses a single font and layout; we have not isolated the effect of typographic choices on the rendered-image gap.

\section{Conclusion}

We present a systematic study of modality effects in VLMs across two evaluation protocols and eight benchmarks. Our findings show that modality dominance is task-conditional: when visual input is redundant (Protocol~A), VLMs consistently default to text, with 96--100\% text preference in the mismatch condition; when visual input is informative (Protocol~B), VLMs correctly leverage it, with image advantages of 10--54pp. This suggests that current VLMs are not fundamentally text-first, but that their apparent text bias reflects rational use of the most informative modality. Future work should investigate the mismatch condition under Protocol~B settings, and whether prompt strategies can reduce unnecessary text dominance in rendered-text scenarios.

\bibliographystyle{acl_natbib}
\bibliography{custom}

\end{document}
